\documentclass[11pt,a4paper]{article}

% Packages
\usepackage[margin=1in]{geometry}
\usepackage{booktabs}
\usepackage{longtable}
\usepackage{enumitem}
\usepackage{xcolor}
\usepackage{hyperref}
\usepackage{listings}
\usepackage{fancyhdr}
\usepackage{titlesec}
\usepackage{amsmath}
\usepackage{graphicx}
\usepackage{microtype}

% Colors
\definecolor{codebg}{RGB}{245,245,245}
\definecolor{codeframe}{RGB}{200,200,200}
\definecolor{linkblue}{RGB}{0,80,160}

% Hyperref config
\hypersetup{
    colorlinks=true,
    linkcolor=linkblue,
    urlcolor=linkblue,
    citecolor=linkblue,
}

% Code listing style
\lstset{
    backgroundcolor=\color{codebg},
    frame=single,
    rulecolor=\color{codeframe},
    basicstyle=\ttfamily\small,
    keywordstyle=\bfseries\color{blue!70!black},
    commentstyle=\color{green!50!black},
    stringstyle=\color{red!60!black},
    breaklines=true,
    breakatwhitespace=true,
    tabsize=2,
    showstringspaces=false,
    language=Python,
    numbers=none,
}

% Header/Footer
\pagestyle{fancy}
\fancyhf{}
\fancyhead[L]{\textsc{svVascularize Optimization Report}}
\fancyhead[R]{\textsc{March 2026}}
\fancyfoot[C]{\thepage}

% Title formatting
\titleformat{\section}{\Large\bfseries\color{black}}{}{0em}{}[\titlerule]
\titleformat{\subsection}{\large\bfseries}{}{0em}{}
\titleformat{\subsubsection}{\normalsize\bfseries}{}{0em}{}

\begin{document}

% Title page
\begin{titlepage}
\centering
\vspace*{3cm}
{\Huge\bfseries svVascularize\\[0.3em] Performance Optimization\\[0.3em] \& Visualization Report\par}
\vspace{1.5cm}
{\Large Commit \texttt{c2001be} on branch \texttt{MAC\_fixed\_protocols}\par}
\vspace{1cm}
{\large March 9, 2026\par}
\vspace{2cm}

\begin{tabular}{rl}
\textbf{Files changed:} & 18 (17 modified + 1 new) \\
\textbf{Lines added:} & +1,005 \\
\textbf{Lines removed:} & $-$934 \\
\textbf{Net change:} & +71 lines \\
\end{tabular}

\vspace{2cm}

{\large\textsc{Categories of Changes}}\\[0.5em]
\begin{tabular}{ll}
\toprule
Category & Technique \\
\midrule
Forest connection pipeline & Vectorization, caching, deduplication \\
Import/meshing pipeline & Algorithm substitution, short-circuit \\
Visualization & Batch rendering, offscreen VTK \\
Code cleanup & Dead code removal, deepcopy elimination \\
\bottomrule
\end{tabular}

\vfill
\end{titlepage}

\tableofcontents
\newpage

%=============================================================================
\section{Executive Summary}
%=============================================================================

This commit delivers a comprehensive set of performance optimizations and visualization fixes to the \texttt{svVascularize} (\texttt{svv}) codebase. The changes target three areas:

\begin{enumerate}[leftmargin=2em]
    \item \textbf{Forest connection pipeline} --- The primary computational bottleneck. Optimizations include Dijkstra result caching, assignment pair deduplication, shared constraint evaluation caching, and full NumPy vectorization of curve evaluation and segment distance computation.
    \item \textbf{Import/meshing pipeline} --- Elimination of redundant computation when importing pre-built meshes, replacement of expensive algorithms with faster alternatives, and in-process TetGen calls.
    \item \textbf{Visualization} --- Batch cylinder rendering replaces per-vessel actor creation, and macOS-specific offscreen VTK rendering fixes resolve GUI hangs on Apple Silicon.
\end{enumerate}

Across all changes, \texttt{deepcopy} was systematically replaced with appropriate shallow-copy alternatives, and approximately 300 lines of dead/commented-out code were removed.

%=============================================================================
\section{Forest Connection Pipeline Optimizations}
%=============================================================================

The forest connection pipeline (\texttt{forest.connect()}) orchestrates the geometric connection of interpenetrating vascular networks. It was identified as the primary bottleneck in end-to-end vascular generation workflows.

%-----------------------------------------------------------------------------
\subsection{Geodesic Caching (\texttt{geodesic.py})}
%-----------------------------------------------------------------------------

\subsubsection{Edge Extraction Vectorization}

\textbf{Before:} Triple-nested Python loop over tetrahedra and 6 edge pairs, checking a Python \texttt{set} for duplicate edges:
\begin{lstlisting}
for tet in tetrahedra:
    for i, j in edge_pairs:
        edge = (tet[i], tet[j])
        if edge not in seen:
            edges.append(edge)
            seen.add(edge)
\end{lstlisting}

\textbf{After:} Single vectorized operation using advanced indexing and \texttt{np.unique}:
\begin{lstlisting}
all_edges = tetra[:, edge_pairs].reshape(-1, 2)
sorted_edges = np.sort(all_edges, axis=1)
unique_edges = np.unique(sorted_edges, axis=0)
lengths = np.linalg.norm(
    pts[unique_edges[:, 0]] - pts[unique_edges[:, 1]], axis=1)
\end{lstlisting}

\textbf{Impact:} Eliminates $O(n_{\text{tetra}} \times 6)$ Python-level loop iterations and set lookups.

\subsubsection{Dijkstra Result Caching}

\textbf{Before:} Every call to \texttt{get\_path(start, end)} ran a full Dijkstra computation from scratch.

\textbf{After:} A dictionary cache \texttt{\_dijkstra\_cache} stores \texttt{(dist, pred)} results keyed by source node. Since Dijkstra from a single source yields shortest paths to \emph{all} nodes, subsequent queries from the same source are $O(n)$ path reconstruction instead of $O(E \log V)$ graph traversal.

\textbf{Impact:} In the forest connection pipeline, many terminal pairs share source nodes. This can reduce Dijkstra computations by 50--80\% depending on topology.

%-----------------------------------------------------------------------------
\subsection{Assignment Deduplication (\texttt{assign.py})}
%-----------------------------------------------------------------------------

\textbf{Problem:} Bidirectional $k$-NN queries between terminal nodes produce duplicate $(i,j)$ pairs when both the forward and reverse searches match the same pair.

\textbf{Before:} All pairs were processed independently, computing geodesic paths for duplicates.

\textbf{After:} \texttt{numpy.unique(pairs, axis=0, return\_inverse=True)} deduplicates before geodesic computation. Results are mapped back via the inverse index.
\begin{lstlisting}
unique_pairs, inverse_idx = np.unique(
    pairs, axis=0, return_inverse=True)
# Compute geodesic only for unique_pairs
for i, (a, b) in enumerate(unique_pairs):
    paths[i] = compute_geodesic(a, b)
# Map back
all_paths = [paths[inv] for inv in inverse_idx]
\end{lstlisting}

\textbf{Additional cleanup:}
\begin{itemize}[nosep]
    \item Extracted helper functions: \texttt{\_make\_linear\_interp()}, \texttt{\_make\_geodesic\_interp()}, \texttt{\_compute\_path\_and\_func()} --- eliminating 6 copy-pasted code blocks.
    \item Replaced string key \texttt{str(i)+','+str(j)} with tuple key \texttt{(int(i), int(j))} for sparse matrix indexing.
    \item Removed ${\sim}200$ lines of commented-out dead code.
\end{itemize}

\textbf{Impact:} Roughly halves the number of geodesic computations; major code simplification ($-$461 lines, $+$193 lines).

%-----------------------------------------------------------------------------
\subsection{Constraint Sharing (\texttt{base\_connection.py})}
%-----------------------------------------------------------------------------

\textbf{Problem:} The SLSQP optimizer evaluates 4 constraint functions per iteration. Each independently rebuilt control points and the curve object --- 75\% redundant work.

\textbf{Before:}
\begin{lstlisting}
def curvature_constraint(x):
    ctrl_pts = self._build_control_points(x)
    curve = Curve(ctrl_pts, ...)  # Rebuilt each time
    ...
def boundary_constraint(x):
    ctrl_pts = self._build_control_points(x)  # Same x!
    curve = Curve(ctrl_pts, ...)  # Rebuilt again!
    ...
\end{lstlisting}

\textbf{After:} A shared \texttt{\_cache} dictionary stores the curve and evaluated points, keyed by the byte representation of the input:
\begin{lstlisting}
_cache = {'key': None, 'ctrl_pts': None, 'curve': None, ...}

def _get_cached(ctrlpts_flat):
    key = ctrlpts_flat.tobytes()
    if key != _cache['key']:
        _cache['key'] = key
        _cache['ctrl_pts'] = self._build_control_points(...)
        _cache['curve'] = Curve(_cache['ctrl_pts'], ...)
    return _cache
\end{lstlisting}

\textbf{Additional optimizations:}
\begin{itemize}[nosep]
    \item Pre-computed constants outside closures: \texttt{min\_radius\_needed}, \texttt{max\_radius}, all \texttt{t\_values} arrays, \texttt{\_collision\_radii}.
    \item Reduced sampling: curvature $100 \to 50$ points; boundary $100 \to 40$ points.
    \item \texttt{np.zeros} replaced with \texttt{np.empty} for segment arrays.
\end{itemize}

\textbf{Impact:} Curve construction cost reduced by $4\times$ per optimizer iteration. Sampling reduction provides additional ${\sim}2\times$ speedup in constraint evaluation.

%-----------------------------------------------------------------------------
\subsection{Bezier Vectorization (\texttt{bezier.py})}
%-----------------------------------------------------------------------------

\textbf{Before:} Per-point Python loop calling scalar De Casteljau:
\begin{lstlisting}
result = np.array([self._de_casteljau(self.control_points, t)
                   for t in t_values])
\end{lstlisting}

\textbf{After:} Vectorized De Casteljau evaluating all $t$ values simultaneously:
\begin{lstlisting}
pts = np.broadcast_to(
    self.control_points[:, None, :],
    (n, len(t_values), d)).copy()
t = t_values[None, :, None]
for i in range(n - 1):
    pts[:n-1-i] = (1 - t) * pts[:n-1-i] + t * pts[1:n-i]
return pts[0]
\end{lstlisting}

The same vectorization was applied to the \texttt{derivative()} method. For the special case of $m=1$ (constant derivative), \texttt{np.broadcast\_to} with \texttt{.copy()} is used.

\textbf{Impact:} Eliminates per-point Python overhead in the inner curve evaluation loop.

%-----------------------------------------------------------------------------
\subsection{Catmull--Rom Vectorization (\texttt{catmullrom.py})}
%-----------------------------------------------------------------------------

\textbf{Before:} Per-point loop: clamp $t$, find segment, evaluate scalar polynomial.

\textbf{After:} Segment-batched vectorization:
\begin{lstlisting}
seg_idx = np.clip(np.floor(t_arr * (n - 1)).astype(int),
                  0, n - 2)
for seg in np.unique(seg_idx):
    mask = (seg_idx == seg)
    u_vals = local_t[mask]
    u2, u3 = u_vals**2, u_vals**3
    # Catmull-Rom polynomial via np.outer
    out[mask] = 0.5 * (2*p1 + np.outer(u_vals, a)
                + np.outer(u2, b) + np.outer(u3, c))
\end{lstlisting}

\textbf{Impact:} Batched evaluation per segment; most $t$ values cluster in a few segments, making this highly efficient.

%-----------------------------------------------------------------------------
\subsection{Segment Distance Vectorization (\texttt{c\_distance.py})}
%-----------------------------------------------------------------------------

\textbf{Problem:} \texttt{minimum\_segment\_distance()} is called inside the SLSQP optimizer's collision constraint --- one of the most frequently invoked functions.

\textbf{Before:} Double Python for-loop computing each segment pair serially with scalar dot products.

\textbf{After:} Full NumPy broadcasting:
\begin{lstlisting}
# Extract endpoints: A0, A1 shape (N,3); B0, B1 shape (M,3)
# Broadcast to (N,1,3) vs (1,M,3) for pairwise computation
d = A0[:,None,:] - B0[None,:,:]
# Compute all dot products, parametric values, distances in bulk
denom = dotuv * dotuv - dotuu * dotvv  # (N, M)
\end{lstlisting}

Falls back to scalar computation only for degenerate/parallel cases where $|\text{denom}| < 10^{-14}$.

A new \texttt{\_point\_to\_segment\_distance\_batch()} helper handles vectorized point-to-segment distance.

\texttt{minimum\_self\_segment\_distance()} now computes the full pairwise matrix in one call, then masks out self-distances and adjacent segments.

\textbf{Impact:} $O(N \times M)$ scalar operations replaced with bulk NumPy; dramatic reduction in Python overhead.

%-----------------------------------------------------------------------------
\subsection{\texttt{deepcopy} Removal}
%-----------------------------------------------------------------------------

Systematic replacement of \texttt{copy.deepcopy()} with appropriate alternatives across 5 files:

\begin{longtable}{p{3.8cm}p{4cm}p{4.2cm}}
\toprule
\textbf{File} & \textbf{Before} & \textbf{After} \\
\midrule
\endhead
\texttt{bifurcation.py} & \texttt{deepcopy(connectivity)} & \texttt{connectivity.copy()} \\
 & \texttt{deepcopy(sorted(set(...)))} & \texttt{sorted(set(...))} \\
 & \texttt{deepcopy(upstream)} & \texttt{list(upstream)} \\
\midrule
\texttt{tree.py} & \texttt{deepcopy(vessel\_map[k])} & Manual dict + \texttt{list()} \\
\midrule
\texttt{vessel\_connection.py} & \texttt{deepcopy(tmp)} on arrays & \texttt{np.empty()} (fresh alloc) \\
\midrule
\texttt{assign.py} & \texttt{deepcopy(spline\_data)} & Fresh allocation \\
\bottomrule
\end{longtable}

\textbf{Rationale:} \texttt{deepcopy} recursively copies every nested object, which is extremely expensive for NumPy arrays (which support efficient \texttt{.copy()}) and for containers of immutable objects (integers, strings) where shallow copies suffice.

%=============================================================================
\section{Import/Meshing Pipeline Optimizations}
%=============================================================================

%-----------------------------------------------------------------------------
\subsection{Skip \texttt{contour()} for Imported Meshes (\texttt{domain.py})}
%-----------------------------------------------------------------------------

\textbf{Problem:} When importing pre-built meshes (STL/VTP), the \texttt{get\_boundary()} method still called \texttt{contour()}, which evaluates the implicit function at $\text{resolution}^3$ grid points to produce \texttt{self.grid}. This grid is never used downstream when an \texttt{original\_boundary} exists.

\textbf{Fix:} When \texttt{original\_boundary} is present, set \texttt{self.grid = None} and skip the \texttt{contour()} call entirely.

\textbf{Impact:} For a resolution of 100, this skips $10^6$ implicit function evaluations.

%-----------------------------------------------------------------------------
\subsection{In-Process TetGen (\texttt{tetrahedralize.py})}
%-----------------------------------------------------------------------------

\textbf{Before:} Always spawned a subprocess: save VTP to disk $\to$ spawn Python process $\to$ run TetGen $\to$ save NPZ $\to$ load NPZ.

\textbf{After:} Calls the \texttt{tetgen} Python binding directly in-process via \texttt{\_run\_tetgen(surface)}. Falls back to the subprocess approach only if the in-process call raises an exception (crash isolation for edge cases where TetGen segfaults).

\textbf{Impact:} Eliminates file I/O and process spawn overhead, which dominated for small-to-medium meshes.

%-----------------------------------------------------------------------------
\subsection{ConvexHull Replaces \texttt{delaunay\_3d} (\texttt{domain.py})}
%-----------------------------------------------------------------------------

\textbf{Before:} \texttt{pv.PolyData(points).delaunay\_3d()} for computing convex hull volume.

\textbf{After:} \texttt{scipy.spatial.ConvexHull(points).volume} --- a purpose-built algorithm that is orders of magnitude faster for this specific task.

%-----------------------------------------------------------------------------
\subsection{Single cKDTree Replaces BallTree (\texttt{domain.py}, \texttt{dmn.py})}
%-----------------------------------------------------------------------------

\textbf{Before:} Both a \texttt{scipy.spatial.cKDTree} and a \texttt{sklearn.neighbors.BallTree} were built on the same cell centers. BallTree was used for \texttt{query\_radius()} calls.

\textbf{After:} \texttt{mesh\_tree\_2 = mesh\_tree} --- a single cKDTree using \texttt{query\_ball\_point()} instead of BallTree's \texttt{query\_radius()}.

\textbf{Impact:} Eliminates the sklearn dependency for this code path and avoids constructing a redundant spatial index.

%-----------------------------------------------------------------------------
\subsection{Vectorized \texttt{evaluate\_fast} (\texttt{domain.py})}
%-----------------------------------------------------------------------------

\textbf{Before:} Python for-loop filling the index array row by row:
\begin{lstlisting}
for i in range(len(indices)):
    inds[i, :len(indices[i])] = indices[i]
\end{lstlisting}

\textbf{After:} Two-phase vectorized fill:
\begin{enumerate}[nosep]
    \item Batch handle empty rows with a boolean mask.
    \item Group rows by identical length using \texttt{np.unique(lengths)}, then fill each group as a single \texttt{np.array} block.
\end{enumerate}

Auto-chunking (\texttt{\_CHUNK=2000}) caps peak memory for large inputs.

%-----------------------------------------------------------------------------
\subsection{Cached \texttt{normalize\_scale} (\texttt{domain.py})}
%-----------------------------------------------------------------------------

Pre-computed once in \texttt{build()} as \texttt{self.\_normalize\_scale} instead of recomputing $\|\max(\text{pts}) - \min(\text{pts})\|$ on every call to \texttt{evaluate\_fast()}.

%-----------------------------------------------------------------------------
\subsection{Thread-Parallel Patch Operations (\texttt{domain.py})}
%-----------------------------------------------------------------------------

\texttt{create()} and \texttt{solve()} now use \texttt{ThreadPoolExecutor} to process patches in parallel, since each patch's creation and solving is independent.

%=============================================================================
\section{Visualization Improvements}
%=============================================================================

%-----------------------------------------------------------------------------
\subsection{Batch Cylinder Rendering (\texttt{batch\_cylinders.py}, NEW)}
%-----------------------------------------------------------------------------

A new module provides vectorized construction of cylinder geometry for efficient VTK rendering.

\textbf{Key function:} \texttt{make\_cylinders\_batch(centers, directions, radii, heights, resolution=8)} constructs all cylinder geometry in NumPy:
\begin{itemize}[nosep]
    \item Builds local coordinate frames $(u, v, w)$ for all cylinders simultaneously.
    \item Generates ring vertices via vectorized trigonometry.
    \item Constructs face connectivity with a precomputed pattern tiled across all cylinders.
    \item Returns a single \texttt{pv.PolyData} object.
\end{itemize}

\textbf{Helper functions:}
\begin{itemize}[nosep]
    \item \texttt{tree\_to\_merged\_mesh(tree)} --- converts a Tree's vessel data into a single merged mesh.
    \item \texttt{segments\_to\_merged\_mesh(segments)} --- converts connection vessel segments into a single merged mesh.
\end{itemize}

\textbf{Impact:} A forest with 1000+ vessels previously created 1000+ VTK actors. Now all vessels are merged into 1--3 actors, dramatically reducing GPU draw calls and VTK overhead.

%-----------------------------------------------------------------------------
\subsection{Adoption of Batch Rendering}
%-----------------------------------------------------------------------------

Three visualization files were updated to use the new batch module:

\begin{longtable}{p{4.5cm}p{4cm}p{3.5cm}}
\toprule
\textbf{File} & \textbf{Before} & \textbf{After} \\
\midrule
\endhead
\texttt{visualize/tree/show.py} & \texttt{trange} loop + per-vessel \texttt{pv.Cylinder()} & Single \texttt{tree\_to\_merged\_mesh()} \\
\texttt{visualize/forest/show.py} & Per-vessel loop + per-connection loop & \texttt{tree\_to\_merged\_mesh()} + \texttt{segments\_to\_merged\_mesh()} \\
\texttt{visualize/gui/vtk\_widget.py} & Per-vessel loop with periodic \texttt{processEvents()} & Single batch call per group \\
\bottomrule
\end{longtable}

%-----------------------------------------------------------------------------
\subsection{macOS VTK/Qt Fixes (\texttt{vtk\_widget.py})}
%-----------------------------------------------------------------------------

The VTK/Qt visualization had critical issues on macOS Apple Silicon with conda environments:
\begin{itemize}[nosep]
    \item \textbf{Offscreen rendering:} Uses offscreen VTK rendering with blit-to-QLabel approach, bypassing \texttt{QtInteractor} deadlock that caused GUI hangs.
    \item \textbf{Software ray-cast picking:} Implements software ray-cast picking for offscreen mode, since VTK's hardware picker segfaults without an on-screen OpenGL context.
\end{itemize}

%=============================================================================
\section{Code Cleanup}
%=============================================================================

%-----------------------------------------------------------------------------
\subsection{Dead Code Removal}
%-----------------------------------------------------------------------------

\begin{itemize}[nosep]
    \item \texttt{assign.py}: Removed ${\sim}200$ lines of commented-out code blocks.
    \item \texttt{base\_connection.py}: Removed ${\sim}30$ lines of commented-out code.
    \item \texttt{tree\_connection.py}: Removed debug \texttt{print()} statements.
\end{itemize}

%-----------------------------------------------------------------------------
\subsection{Helper Function Extraction}
%-----------------------------------------------------------------------------

\begin{itemize}[nosep]
    \item \texttt{assign.py}: \texttt{\_make\_linear\_interp()}, \texttt{\_make\_geodesic\_interp()}, \texttt{\_compute\_path\_and\_func()} --- replacing 6 copy-pasted blocks.
    \item \texttt{vessel\_connection.py}: \texttt{\_build\_tree\_collision\_segments()} --- extracted from inline index manipulation code.
\end{itemize}

%=============================================================================
\section{Summary of Changes by File}
%=============================================================================

\begin{longtable}{p{5cm}ccp{4.5cm}}
\toprule
\textbf{File} & \textbf{+Lines} & \textbf{$-$Lines} & \textbf{Primary Change} \\
\midrule
\endhead
\texttt{CLAUDE.md} & 29 & 0 & Documentation \\
\texttt{domain/domain.py} & 116 & 72 & ConvexHull, skip contour, vectorize evaluate\_fast, threading \\
\texttt{domain/io/dmn.py} & 4 & 10 & cKDTree reuse \\
\texttt{domain/routines/tetrahedralize.py} & 98 & 83 & In-process TetGen \\
\texttt{forest/connect/assign.py} & 193 & 461 & Pair deduplication, helpers \\
\texttt{forest/connect/base\_connection.py} & 108 & 146 & Constraint caching \\
\texttt{forest/connect/bezier.py} & 22 & 4 & Vectorized De Casteljau \\
\texttt{forest/connect/catmullrom.py} & 33 & 32 & Segment-batched evaluation \\
\texttt{forest/connect/geodesic.py} & 57 & 45 & Dijkstra cache, vectorized edges \\
\texttt{forest/connect/tree\_connection.py} & 9 & 26 & Cleanup, remove debug prints \\
\texttt{forest/connect/vessel\_connection.py} & 34 & 73 & deepcopy removal, helper extraction \\
\texttt{tree/branch/bifurcation.py} & 13 & 13 & deepcopy $\to$ shallow copy \\
\texttt{tree/tree.py} & 2 & 1 & deepcopy $\to$ manual dict copy \\
\texttt{utils/spatial/c\_distance.py} & 139 & 47 & Full NumPy vectorization \\
\texttt{visualize/batch\_cylinders.py} & 239 & 0 & New: batch cylinder module \\
\texttt{visualize/forest/show.py} & 14 & 28 & Batch rendering adoption \\
\texttt{visualize/gui/vtk\_widget.py} & 12 & 48 & Batch rendering + offscreen fixes \\
\texttt{visualize/tree/show.py} & 4 & 14 & Batch rendering adoption \\
\midrule
\textbf{Total} & \textbf{1,005} & \textbf{934} & \textbf{Net: +71 lines} \\
\bottomrule
\end{longtable}

%=============================================================================
\section{Optimization Techniques Summary}
%=============================================================================

\begin{longtable}{p{3.5cm}p{3.5cm}p{5cm}}
\toprule
\textbf{Technique} & \textbf{Files} & \textbf{Description} \\
\midrule
\endhead
NumPy Vectorization & bezier, catmullrom, c\_distance, geodesic, domain & Replace Python loops with broadcasting and advanced indexing \\
\midrule
Caching / Memoization & base\_connection, geodesic, domain & Closure-level cache, Dijkstra cache, normalize\_scale \\
\midrule
\texttt{deepcopy} Removal & bifurcation, tree, vessel\_connection, assign & \texttt{.copy()}, \texttt{list()}, fresh allocation \\
\midrule
Deduplication & assign & \texttt{np.unique} on $(i,j)$ pairs before geodesic \\
\midrule
Algorithm Substitution & domain, tetrahedralize & ConvexHull for delaunay\_3d; in-process TetGen \\
\midrule
Dependency Elimination & domain, dmn & Removed sklearn BallTree; reuse cKDTree \\
\midrule
Short-Circuit & domain & Skip \texttt{contour()} for imported meshes \\
\midrule
Batch Rendering & batch\_cylinders, vtk\_widget, forest/show, tree/show & Single merged PolyData per group \\
\midrule
Thread Parallelism & domain & ThreadPoolExecutor for patch creation/solving \\
\bottomrule
\end{longtable}

\end{document}
